\begin{exsol}{18}{\MR{} Food for Thought 5.2}{164}{Exam}{B2}
\asked A ``signal of intent'' (e.g.\ a price cut announcing that an incumbent will fight an
expansion) is not credible if, once the rival has sunk its costs, fighting no longer pays. What
means can make such signals credible?

\method Commitment and reneging (\MR{} p.~133), the box on terrorism and parenting (p.~135),
noncontractual commitment and its costs (p.~139), reputation (pp.~139--140), signaling
(pp.~154--156), and the logic of subgame perfection: a threat counts only if carrying it out is
optimal in the subgame where it must be carried out. Commitment by removing options is also in
Cabral's \S7.2 and Varian's Ch.~30 (``Games of Commitment'').

\solution
\textbf{The problem in one inequality.} Let the incumbent earn $M$ if the rival stays out, and
after entry $F$ if it fights and $A$ if it accommodates, with $M > A > F$. Fighting is then not
optimal after entry, so the rival, reasoning backwards, ignores any announcement. A price cut
today is ``cheap'' in the relevant sense: it is reversible and leaves $F < A$ untouched. A signal
becomes credible only if it (i) changes the post-entry payoffs so that $F \ge A$, (ii) removes
the option to accommodate, (iii) links this market to others or to the future, or (iv) reveals
information that the incumbent really is a fighter. Each route has a cost, and the device is worth
buying only if the deterrence gain $M - A$ exceeds it.

\textbf{1. Sunk, specific investments.} Build capacity, distribution or a plant usable only in
this market. Sunk capacity lowers the marginal cost of producing a lot, so aggressive output
becomes the best reply after entry. Apple's deliberately overspecialized Macintosh plant,
heavily publicized, made dropping the product costly and so committed Apple to the market in the
view of its rivals, employees and customers (p.~133). \emph{Cost:} the investment itself, a plant ``that may become obsolete
sooner'' (p.~139), and the loss if entry happens anyway.

\textbf{2. Burning bridges.} Destroy the option to retreat: William the Conqueror's and
Cort\'es's fleets (p.~133), or a public, irreversible exit from the alternatives. \emph{Cost:}
the forgone option has value in the states where retreat would have been right; the fleets were
``perfectly good'' (p.~139).

\textbf{3. Contracts with third parties.} Write the threat into obligations to someone else, so
that accommodating triggers a penalty. Examples: price-matching or most-favoured-customer clauses
that make a price war automatic and price increases costly; supply commitments; bonds or public
pledges with enforceable penalties; licensing a third producer so that a promised price cannot be
raised later (the logic of \MR's Analytical Problem 5.6, p.~165). \emph{Weakness:} renegotiation.
If the third party gains by releasing the incumbent ex post, it will (pp.~134--135), so the
contract must give the third party a reason to insist. \emph{Cost:} lost flexibility and,
sometimes, antitrust exposure.

\textbf{4. Delegation to an agent with different preferences.} Put the post-entry decision in the
hands of someone who \emph{wants} to fight: a manager known to be aggressive, or a division paid on
market share or sales rather than profit. The agent's best reply after entry is to fight, which
is what deters. \emph{Cost:} the agent also fights when it should not and distorts other
decisions; this is an agency cost paid in every state.

\textbf{5. Reputation.} If the incumbent faces many potential entrants, markets or periods,
accommodating once invites further challenges. Concern for future effects ``can stiffen resolve''
(box, p.~135): fighting today protects profits elsewhere, so the comparison becomes $F$ plus
the value of reputation against $A$. Its strength grows with the frequency of similar
encounters, the horizon and the profitability at stake (p.~139). \emph{Cost:} the incumbent must
actually fight sometimes, and the device weakens near the end of a known horizon.

\textbf{6. Costly signals that reveal type.} If the doubt concerns the incumbent's costs or
toughness (private information), a deep price cut that only a low-cost or committed firm can
afford separates types, as in Spence signaling (pp.~154--156). Credibility comes from the
signal being too expensive for the weak type to mimic. \emph{Cost:} profits sacrificed to
signaling, which is wasteful from the joint point of view.

\textbf{Common thread.} All six turn talk into a changed game: they alter payoffs, options or
information so that the announced action is subgame perfect. Pure announcements, however
emphatic, do not.

\interp Credibility is a property of incentives after the fact, not of words before it. For
managers this means two things: commitment has a price, so the value of deterrence must be weighed
against lost flexibility; and devices that can be renegotiated or mimicked are weak. The same
reasoning explains why delegation and organizational design, such as who decides and on what they
are paid, are strategic choices and not only internal ones.

\verify Consistency checks: each device makes the announced action optimal in the post-entry
subgame (backward induction), which is the test used for the hostage and parenting examples on
p.~135 and for Cort\'es in Cabral's Exercise 7.6 (item 10). Every source citation was checked
against \MR's text (pp.~133--140, 154--156, 165).

\answer{A signal of intent is credible only if the signaled action is optimal when the time
comes. The means are: sunk, market-specific investments (capacity, a specialized plant);
burning bridges; contracts with third parties that penalize backing down and resist
renegotiation; delegation to an agent with more aggressive preferences or incentives; reputation
across markets or periods; and costly signals that a weak type cannot afford. Each costs
flexibility or resources and is worth using only if deterrence gains exceed that cost.}
\end{exsol}
